% problem_id: S001-more-groupoids-lemma-slice
% source_id: S001 (Stacks Project)
% source_file: more-groupoids.tex
% statement_locator: more-groupoids.tex lines 1849--1875
% proof_locator: more-groupoids.tex lines 1877--2019
% env: lemma
% id_note: label 跨文件重名 ⇒ id 用文件限定形式
% pairing: tight (verified)
% status: complete (statement + author proof verbatim + reason + locators)
%
\begin{lemma}
\label{lemma-slice}
Let $S$ be a scheme.
Let $(U, R, s, t, c, e, i)$ be a groupoid scheme over $S$.
Let $G \to U$ be the stabilizer group scheme.
Assume $s$ and $t$ are Cohen-Macaulay and locally of finite presentation.
Let $u \in U$ be a finite type point of the scheme $U$, see
Morphisms, Definition \ref{morphisms-definition-finite-type-point}.
With notation as in
Situation \ref{situation-slice},
set
$$
d_1 = \dim(G_u), \quad
d_2 = \dim_{e(u)}(F_u).
$$
If $d_2 > d_1$, then there exist an affine scheme $U'$
and a morphism $g : U' \to U$ such that (with notation as in
Situation \ref{situation-slice})
\begin{enumerate}
\item $g$ is an immersion
\item $u \in U'$,
\item $g$ is locally of finite presentation,
\item the morphism $h : U' \times_{g, U, t} R \longrightarrow U$
is Cohen-Macaulay at $(u, e(u))$, and
\item we have $\dim_{e'(u)}(F'_u) = d_2 - 1$.
\end{enumerate}
\end{lemma}

\begin{proof}
Let $\Spec(A) \subset U$ be an affine neighbourhood of $u$
such that $u$ corresponds to a closed point of $U$, see
Morphisms, Lemma \ref{morphisms-lemma-identify-finite-type-points}.
Let $\Spec(B) \subset R$ be an affine neighbourhood of $e(u)$
which maps via $j$ into the open
$\Spec(A) \times_S \Spec(A) \subset U \times_S U$.
Let $\mathfrak m \subset A$ be the maximal ideal corresponding to $u$.
Let $\mathfrak q \subset B$ be the prime ideal corresponding to $e(u)$.
Pictures:
$$
\vcenter{
\xymatrix{
B & A \ar[l]^s \\
A \ar[u]^t
}
}
\quad\text{and}\quad
\vcenter{
\xymatrix{
B_{\mathfrak q} & A_{\mathfrak m} \ar[l]^s \\
A_{\mathfrak m} \ar[u]^t
}
}
$$
Note that the two induced maps
$s, t : \kappa(\mathfrak m) \to \kappa(\mathfrak q)$
are equal and isomorphisms as $s \circ e = t \circ e = \text{id}_U$.
In particular we see that $\mathfrak q$
is a maximal ideal as well. The ring maps $s, t : A \to B$ are
of finite presentation and flat. By assumption the ring
$$
\mathcal{O}_{F_u, e(u)} = B_{\mathfrak q}/s(\mathfrak m)B_{\mathfrak q}
$$
is Cohen-Macaulay of dimension $d_2$. The equality of dimension holds by
Morphisms, Lemma \ref{morphisms-lemma-dimension-fibre-at-a-point}.

\medskip\noindent
Let $R''$ be the restriction of $R$ to $u = \Spec(\kappa(u))$
via the morphism $\Spec(\kappa(u)) \to U$.
As $u \to U$ is locally of finite type,
we see that $(\Spec(\kappa(u)), R'', s'', t'', c'')$
is a groupoid scheme with $s'', t''$ locally of finite type, see
Lemma \ref{lemma-restrict-preserves-type}.
By
Lemma \ref{lemma-groupoid-on-field-dimension-equal-stabilizer}
this implies that $\dim(G'') = \dim(R'')$. We also have
$\dim(R'') = \dim_{e''}(R'') = \dim(\mathcal{O}_{R'', e''})$, see
Lemma \ref{lemma-groupoid-on-field-locally-finite-type-dimension}.
By
Groupoids, Lemma \ref{groupoids-lemma-restrict-stabilizer}
we have $G'' = G_u$. Hence we conclude that
$\dim(\mathcal{O}_{R'', e''}) = d_1$.

\medskip\noindent
As a scheme $R''$ is
$$
R'' =
R \times_{(U \times_S U)}
\Big(
\Spec(\kappa(\mathfrak m)) \times_S \Spec(\kappa(\mathfrak m))
\Big)
$$
Hence an affine open neighbourhood of $e''$ is the spectrum of the ring
$$
B \otimes_{(A \otimes A)} (\kappa(\mathfrak m) \otimes \kappa(\mathfrak m))
=
B/s(\mathfrak m)B + t(\mathfrak m)B
$$
We conclude that
$$
\mathcal{O}_{R'', e''} =
B_{\mathfrak q}/s(\mathfrak m)B_{\mathfrak q} + t(\mathfrak m)B_{\mathfrak q}
$$
and so now we know that this ring has dimension $d_1$.

\medskip\noindent
We claim this implies we can find
an element $f \in \mathfrak m$ such that
$$
\dim(B_{\mathfrak q}/(s(\mathfrak m)B_{\mathfrak q} + fB_{\mathfrak q}) < d_2
$$
Namely, suppose $\mathfrak n_j \supset s(\mathfrak m)B_{\mathfrak q}$,
$j = 1, \ldots, m$ correspond to the minimal primes of the local ring
$B_{\mathfrak q}/s(\mathfrak m)B_{\mathfrak q}$.
There are finitely many as this ring is Noetherian (since it is essentially
of finite type over a field -- but also because a Cohen-Macaulay ring is
Noetherian). By the Cohen-Macaulay condition we have
$\dim(B_{\mathfrak q}/\mathfrak n_j) = d_2$, for example by
Algebra, Lemma \ref{algebra-lemma-CM-dim-formula}.
Note that
$\dim(B_{\mathfrak q}/(\mathfrak n_j + t(\mathfrak m)B_{\mathfrak q}))
\leq d_1$
as it is a quotient of the ring
$\mathcal{O}_{R'', e''} =
B_{\mathfrak q}/s(\mathfrak m)B_{\mathfrak q} + t(\mathfrak m)B_{\mathfrak q}$
which has dimension $d_1$. As $d_1 < d_2$ this implies that
$\mathfrak m \not \subset t^{-1}(\mathfrak n_i)$.
By prime avoidance, see
Algebra, Lemma \ref{algebra-lemma-silly},
we can find $f \in \mathfrak m$ with $t(f) \not \in \mathfrak n_j$ for
$j = 1, \ldots, m$. For this choice of $f$ we have
the displayed inequality above, see
Algebra, Lemma \ref{algebra-lemma-one-equation}.

\medskip\noindent
Set $A' = A/fA$ and $U' = \Spec(A')$. Then it is clear that
$U' \to U$ is an immersion, locally of finite presentation
and that $u \in U'$. Thus (1), (2) and (3) of the lemma hold.
The morphism
$$
U' \times_{g, U, t} R \longrightarrow U
$$
factors through $\Spec(A)$ and corresponds to the ring map
$$
\xymatrix{
B/t(f)B \ar@{=}[r] & A/(f) \otimes_{A, t} B & A \ar[l]_-s
}
$$
Now, we see $t(f)$ is not a zerodivisor on
$B_{\mathfrak q}/s(\mathfrak m)B_{\mathfrak q}$ as this is a
Cohen-Macaulay ring of positive dimension and $f$ is not contained
in any minimal prime, see for example
Algebra, Lemma \ref{algebra-lemma-reformulate-CM}.
Hence by
Algebra, Lemma \ref{algebra-lemma-grothendieck-general}
we conclude that $s : A_{\mathfrak m} \to B_{\mathfrak q}/t(f)B_{\mathfrak q}$
is flat with fibre ring
$B_{\mathfrak q}/(s(\mathfrak m)B_{\mathfrak q} + t(f)B_{\mathfrak q})$
which is Cohen-Macaulay by
Algebra, Lemma \ref{algebra-lemma-reformulate-CM}
again. This implies part (4) of the lemma.
To see part (5) note that by Diagram (\ref{equation-restriction})
the fibre $F'_u$ is equal to the fibre of $h$ over $u$.
Hence
$\dim_{e'(u)}(F'_u) =
\dim(B_{\mathfrak q}/(s(\mathfrak m)B_{\mathfrak q} + t(f)B_{\mathfrak q}))$
by
Morphisms, Lemma \ref{morphisms-lemma-dimension-fibre-at-a-point}
and the dimension of this ring is $d_2 - 1$ by
Algebra, Lemma \ref{algebra-lemma-reformulate-CM}
once more. This proves the final assertion of the lemma and we win.
\end{proof}
